%% 摘要

% 中文摘要
\begin{abstract}
    边界表示（Boundary Representation, B-rep）是工业 CAD 内核描述实体几何的核心数据结构，既承载 NURBS 曲面、曲线等连续几何，又记录面、边、顶点之间的拓扑引用关系。如何让神经网络直接生成可被几何内核读取、校验并继续编辑的 B-rep，仍是生成式 CAD 研究中的关键问题。现有方法多借助参数化操作序列、离散码本或结构化潜变量作为中间表征，虽然降低了建模难度，但也使几何精度和拓扑闭合质量受到中间表征容量与解码过程的限制。

    针对上述问题，本文提出一种不依赖显式 VQ-VAE 码本的点级 NURBS 序列化方法，并构建端到端 B-rep 自回归生成框架。该方法将 NURBS 曲面与曲线上的采样点、局部面索引和结构控制符统一编码为离散 token 序列，使模型直接在几何采样点和拓扑索引层面学习生成规律。模型采用标准 next-token 自回归 Transformer，对序列中的坐标 token、拓扑索引 token 和结构 token 进行联合建模。为评估生成结果的工程可用性，本文基于 OpenCASCADE 实现了从生成 token 序列到反量化、NURBS 拟合、拓扑缝合、Solid 构造、STEP 写出和 BRepCheck 校验的完整重建流水线，从而将生成质量落实为可复核的几何内核有效性指标。

    实验在 DeepCAD 去重子集上进行，训练集 64\,351 条、验证集 6\,096 条。主模型 hr2\_patch1 训练 200 个 epoch 后，验证集交叉熵收敛至 0.10486，对应困惑度约 1.110。在 1\,005 条独立采样序列上，模型全部成功写出 STEP 文件，其中 980 个实体通过 BRepCheck 校验，of saved 与 of attempts 口径下的有效率均为 97.51\%，明显高于毕业设计任务书提出的 60\% 有效率要求。早期 hr2 基线在 of saved 口径下的有效率为 90.40\%，本文将其作为工程进展对照，而非严格消融实验。实验结果表明，在当前 DeepCAD 子集和无条件生成设定下，点级 NURBS 序列化与标准自回归建模路线能够形成可训练、可重建、可校验的 B-rep 生成闭环；跨数据集泛化能力和具体组件的独立贡献仍需后续实验进一步验证。
    \keywordslist{B-rep, NURBS, CAD 生成, 点级序列化, 自回归 Transformer, 几何重建}
\end{abstract}

% 英文摘要
\begin{engabstract}
    Boundary Representation (B-rep) is the core data structure used by industrial CAD kernels to describe solid geometry, including both continuous NURBS geometry and topological references among faces, edges, and vertices. Directly generating B-rep models that can be read, checked, and further edited by a geometric kernel remains a key challenge in generative CAD. Existing methods often rely on CAD operation sequences, discrete codebooks, or structured latent variables as intermediate representations. These choices ease learning, but they also constrain geometric accuracy and topological closure through the capacity and decoding quality of the intermediate representation.

    To address this bottleneck, this thesis proposes a codebook-free, point-level NURBS serialization method and builds an end-to-end autoregressive framework for B-rep generation. The method encodes sampled NURBS surface and curve points, local face indices, and structural symbols into a unified token sequence, allowing the model to learn directly from geometric samples and topological indices. A standard next-token autoregressive Transformer jointly models coordinate, topology-index, and structure tokens. An OpenCASCADE-based reconstruction pipeline then converts generated sequences into STEP files through dequantization, NURBS fitting, topology sewing, Solid construction, and BRepCheck validation.

    Experiments are conducted on a deduplicated DeepCAD subset with 64\,351 training samples and 6\,096 validation samples. After 200 epochs, the main model hr2\_patch1 reaches a validation cross-entropy of 0.10486, corresponding to a perplexity of about 1.110. On 1\,005 independently sampled sequences, all samples are exported as STEP files and 980 pass BRepCheck, giving a valid rate of 97.51\% under both the of-saved and of-attempts metrics. This result is well above the 60\% requirement of the undergraduate thesis brief. An earlier hr2 baseline achieves 90.40\% under the of-saved metric and is used only as an engineering-progress reference rather than a strict ablation. The results indicate that, under the current DeepCAD subset and unconditional generation setting, point-level NURBS serialization with standard autoregressive modeling can form a trainable, reconstructable, and verifiable B-rep generation pipeline. Cross-dataset generalization and component-level ablations remain future work.
    \engkeywordslist{B-rep, NURBS, CAD generation, point-level serialization, autoregressive Transformer, geometric reconstruction}
\end{engabstract}
